\begin{minipage}{\linewidth}
\vspace{1em}\noindent\textbf{1. When adjusting your mirrors to carry a passenger, the passenger should:}

\begin{itemize}[nosep,leftmargin=2em]
\item [A.] Stand behind the motorcycle.
\item [B.] Sit on the seat with you.
\item [C.] Adjust the mirror while you sit on the motorcycle.
\item [D.] Stand beside the motorcycle in the blind spot.

\end{itemize}
\end{minipage}


\begin{minipage}{\linewidth}
\vspace{1em}\noindent\textbf{2. A passenger should:}

\begin{itemize}[nosep,leftmargin=2em]
\item [A.] Hold onto the seat strap.
\item [B.] Hold onto the side of the seat.
\item [C.] Hold onto the operator's waist.
\item [D.] Not ride without a backrest.

\end{itemize}
\end{minipage}


\begin{minipage}{\linewidth}
\vspace{1em}\noindent\textbf{3. When the motorcycle is stopped at an intersection, a passenger should:}

\begin{itemize}[nosep,leftmargin=2em]
\item [A.] Keep both feet on the foot pegs.
\item [B.] Place one foot on the ground.
\item [C.] Place both feet on the ground.
\item [D.] Move back on the seat.

\end{itemize}
\end{minipage}


\begin{minipage}{\linewidth}
\vspace{1em}\noindent\textbf{4. When carrying a passenger:}

\begin{itemize}[nosep,leftmargin=2em]
\item [A.] You will have more balance.
\item [B.] You will be able to slow more quickly.
\item [C.] It will take longer to accelerate and slow down.
\item [D.] You will have less chance of skidding.

\end{itemize}
\end{minipage}


\begin{minipage}{\linewidth}
\vspace{1em}\noindent\textbf{5. When riding with passengers, you should:}

\begin{itemize}[nosep,leftmargin=2em]
\item [A.] Start slowing sooner.
\item [B.] Use only your front brake to slow.
\item [C.] Never talk to them.
\item [D.] Start in the center of the lane.

\end{itemize}
\end{minipage}


\begin{minipage}{\linewidth}
\vspace{1em}\noindent\textbf{6. When carrying loads:}

\begin{itemize}[nosep,leftmargin=2em]
\item [A.] Keep the load behind the rear axle.
\item [B.] Secure the load with a rope.
\item [C.] Fasten the load against the frame on the back of the seat.
\item [D.] Keep the load low and as far forward as possible.

\end{itemize}
\end{minipage}


\begin{minipage}{\linewidth}
\vspace{1em}\noindent\textbf{7. A group of riders move from a staggered formation into a single file when:}

\begin{itemize}[nosep,leftmargin=2em]
\item [A.] Riding at night.
\item [B.] Traffic is heavy.
\item [C.] Stopping at intersections.
\item [D.] Turning corners.

\end{itemize}
\end{minipage}


\begin{minipage}{\linewidth}
\vspace{1em}\noindent\textbf{8. In this picture, riders are in a staggered formation. The following distance between A and C should be at least:}

\vspace{1ex}
\centering
\includegraphics[width=0.8\linewidth]{ q_8.jpeg }
\captionof*{figure}{ Four motorcycles (labeled A, B, C, D) riding in a staggered formation along a multi-lane road, where A is a leader in left portion, B after leader in right portion, C is third, in left portion and behind the leader. }

\begin{itemize}[nosep,leftmargin=2em]
\item [A.] Four seconds.
\item [B.] One car length.
\item [C.] One second.
\item [D.] Two seconds.

\end{itemize}
\end{minipage}


\begin{minipage}{\linewidth}
\vspace{1em}\noindent\textbf{9. Riders in a staggered formation will be passing a car. After the lead rider passes, he/she should:}

\begin{itemize}[nosep,leftmargin=2em]
\item [A.] Return to the right portion of the lane.
\item [B.] Return to the center portion of the lane.
\item [C.] Return to the left portion of the lane.
\item [D.] Signal the others to follow.

\end{itemize}
\end{minipage}


\begin{minipage}{\linewidth}
\vspace{1em}\noindent\textbf{10. Alcohol enters the bloodstream and quickly reaches the brain. What is affected first?}

\begin{itemize}[nosep,leftmargin=2em]
\item [A.] Judgment.
\item [B.] Speech.
\item [C.] Balance.
\item [D.] Vision.

\end{itemize}
\end{minipage}


\begin{minipage}{\linewidth}
\vspace{1em}\noindent\textbf{11. Alcohol first affects your:}

\begin{itemize}[nosep,leftmargin=2em]
\item [A.] Vision.
\item [B.] Speech.
\item [C.] Judgment.
\item [D.] Balance.

\end{itemize}
\end{minipage}


\begin{minipage}{\linewidth}
\vspace{1em}\noindent\textbf{12. Approximately half of all motorcycle crash fatalities are caused by:}

\begin{itemize}[nosep,leftmargin=2em]
\item [A.] Road conditions.
\item [B.] The rider's helmet.
\item [C.] Alcohol Impairment.
\item [D.] Mechanical Failure.

\end{itemize}
\end{minipage}


\begin{minipage}{\linewidth}
\vspace{1em}\noindent\textbf{13. If you don't control your drinking:}

\begin{itemize}[nosep,leftmargin=2em]
\item [A.] Ride slower.
\item [B.] Allow at least an hour for each drink before riding.
\item [C.] Exercise or dance to wear the drinks off.
\item [D.] Drink hot coffee to be alert

\end{itemize}
\end{minipage}


\begin{minipage}{\linewidth}
\vspace{1em}\noindent\textbf{14. Studies show that almost half of all riders killed in motorcycle crashes:}

\begin{itemize}[nosep,leftmargin=2em]
\item [A.] Had poor eyesight.
\item [B.] Had been drinking.
\item [C.] Were speeding.
\item [D.] Were unskilled.

\end{itemize}
\end{minipage}


\begin{minipage}{\linewidth}
\vspace{1em}\noindent\textbf{15. The ability to think clearly and ride safely is affected by as little as:}

\begin{itemize}[nosep,leftmargin=2em]
\item [A.] 1 drink.
\item [B.] 6 drinks.
\item [C.] 3 drinks.
\item [D.] 2 drinks.

\end{itemize}
\end{minipage}


\begin{minipage}{\linewidth}
\vspace{1em}\noindent\textbf{16. If you drink 4 typical drinks, what is the least amount of alcohol remaining in your body after 2 hrs?}

\begin{itemize}[nosep,leftmargin=2em]
\item [A.] 2 drinks.
\item [B.] 3 drinks.
\item [C.] 1 drink.
\item [D.] 4 drinks.

\end{itemize}
\end{minipage}


\begin{minipage}{\linewidth}
\vspace{1em}\noindent\textbf{17. For greatest safety, your helmet should:}

\begin{itemize}[nosep,leftmargin=2em]
\item [A.] Fit tightly at the base; loosely at the top.
\item [B.] Fit loosely all the way around.
\item [C.] Fit snugly all the way around.
\item [D.] Be worn only on long rides.

\end{itemize}
\end{minipage}


\begin{minipage}{\linewidth}
\vspace{1em}\noindent\textbf{18. The single most important thing you can do to improve your chances of surviving a crash is to:}

\begin{itemize}[nosep,leftmargin=2em]
\item [A.] Lay the motorcycle down before hitting an obstacle.
\item [B.] Wear a securely fastened helmet.
\item [C.] Be familiar with the motorcycle.
\item [D.] Ride in the center portion of the lane.

\end{itemize}
\end{minipage}


\begin{minipage}{\linewidth}
\vspace{1em}\noindent\textbf{19. Wearing a helmet can reduce the chance of a fatal head injury in:}

\begin{itemize}[nosep,leftmargin=2em]
\item [A.] Low-speed crash only.
\item [B.] Single-vehicle crashes only.
\item [C.] Any crash, regardless of speed.
\item [D.] High-speed crash only.

\end{itemize}
\end{minipage}


\begin{minipage}{\linewidth}
\vspace{1em}\noindent\textbf{20. Wearing a helmet will:}

\begin{itemize}[nosep,leftmargin=2em]
\item [A.] Reduce your chances of neck and head injuries.
\item [B.] Increase your chances of a neck injury.
\item [C.] Reduce your vision by two-thirds.
\item [D.] Neither reduce nor increase your chances of a head injury.

\end{itemize}
\end{minipage}


\begin{minipage}{\linewidth}
\vspace{1em}\noindent\textbf{21. Of the following, which one gives the most eye and face protection while riding?}

\begin{itemize}[nosep,leftmargin=2em]
\item [A.] Goggles.
\item [B.] A windshield.
\item [C.] A faceshield.
\item [D.] Sunglasses.

\end{itemize}
\end{minipage}


\begin{minipage}{\linewidth}
\vspace{1em}\noindent\textbf{22. While riding an unfamiliar motorcycle:}

\begin{itemize}[nosep,leftmargin=2em]
\item [A.] Do not shift above third gear.
\item [B.] Stay in the right portion of the lane.
\item [C.] Have the owner ride with you.
\item [D.] Allow extra room for stopping.

\end{itemize}
\end{minipage}


\begin{minipage}{\linewidth}
\vspace{1em}\noindent\textbf{23. The engine cut-off switch is usually located:}

\begin{itemize}[nosep,leftmargin=2em]
\item [A.] Near the clutch.
\item [B.] On the left handlebar.
\item [C.] Near the right hand grip.
\item [D.] On the top of the gas tank.

\end{itemize}
\end{minipage}


\begin{minipage}{\linewidth}
\vspace{1em}\noindent\textbf{24. With a properly adjusted mirror you can see:}

\begin{itemize}[nosep,leftmargin=2em]
\item [A.] The lane beside you.
\item [B.] The lane behind you.
\item [C.] Part of the lane beside you and behind you.
\item [D.] Everything behind and to the side.

\end{itemize}
\end{minipage}


\begin{minipage}{\linewidth}
\vspace{1em}\noindent\textbf{25. When slowing for a normal stop:}

\begin{itemize}[nosep,leftmargin=2em]
\item [A.] Use only the rear brake.
\item [B.] Use less throttle and the front brake only.
\item [C.] Shift down through the gears as you slow.
\item [D.] Pull in the clutch and coast.

\end{itemize}
\end{minipage}


\begin{minipage}{\linewidth}
\vspace{1em}\noindent\textbf{26. When turning, you should:}

\begin{itemize}[nosep,leftmargin=2em]
\item [A.] Watch the center line.
\item [B.] Look straight ahead.
\item [C.] Look at the guard rail.
\item [D.] Look through the turn.

\end{itemize}
\end{minipage}


\begin{minipage}{\linewidth}
\vspace{1em}\noindent\textbf{27. Which one of the following statements is correct:}

\begin{itemize}[nosep,leftmargin=2em]
\item [A.] Press left, lean left, go right.
\item [B.] Press left, lean left, go left.
\item [C.] Press left, lean right, go right.
\item [D.] Press right, lean right, go left.

\end{itemize}
\end{minipage}


\begin{minipage}{\linewidth}
\vspace{1em}\noindent\textbf{28. While in a turn, it is best to:}

\begin{itemize}[nosep,leftmargin=2em]
\item [A.] Maintain a steady speed or gradually accelerate.
\item [B.] Slow down.
\item [C.] Use your brakes.
\item [D.] Accelerate and look straight ahead.

\end{itemize}
\end{minipage}


\begin{minipage}{\linewidth}
\vspace{1em}\noindent\textbf{29. If you need to slow down while turning:}

\begin{itemize}[nosep,leftmargin=2em]
\item [A.] Apply the rear brake, and lean farther into the turn.
\item [B.] Never apply the brakes.
\item [C.] Use caution, apply both brakes.
\item [D.] Apply the rear brake only.

\end{itemize}
\end{minipage}


\begin{minipage}{\linewidth}
\vspace{1em}\noindent\textbf{30. To make a normal stop, use:}

\begin{itemize}[nosep,leftmargin=2em]
\item [A.] The rear brake first.
\item [B.] Both brakes, and downshift.
\item [C.] The front brake and downshift.
\item [D.] The rear brake only.

\end{itemize}
\end{minipage}


\begin{minipage}{\linewidth}
\vspace{1em}\noindent\textbf{31. Your motorcycle has two brakes. Use both brakes:}

\begin{itemize}[nosep,leftmargin=2em]
\item [A.] Only for emergency stops.
\item [B.] While swerving.
\item [C.] Every time you slow or stop.
\item [D.] Only for normal stops.

\end{itemize}
\end{minipage}


\begin{minipage}{\linewidth}
\vspace{1em}\noindent\textbf{32. In a normal turn:}

\begin{itemize}[nosep,leftmargin=2em]
\item [A.] Lean your motorcycle more than your body.
\item [B.] Lean your motorcycle and your body at the same angle.
\item [C.] Do not lean.
\item [D.] Lean your body more than your motorcycle.

\end{itemize}
\end{minipage}


\begin{minipage}{\linewidth}
\vspace{1em}\noindent\textbf{33. In slow tight turns:}

\begin{itemize}[nosep,leftmargin=2em]
\item [A.] Lean the motorcycle only, and keep your body straight.
\item [B.] Lean your body more than the motorcycle.
\item [C.] Lean with the motorcycle.
\item [D.] Do not lean.

\end{itemize}
\end{minipage}


\begin{minipage}{\linewidth}
\vspace{1em}\noindent\textbf{34. To reduce speed before turning:}

\begin{itemize}[nosep,leftmargin=2em]
\item [A.] Shift to neutral.
\item [B.] Pull in the clutch and use both brakes.
\item [C.] Apply the rear brake only.
\item [D.] Close the throttle and apply both brakes.

\end{itemize}
\end{minipage}


\begin{minipage}{\linewidth}
\vspace{1em}\noindent\textbf{35. When turning, it is important to:}

\begin{itemize}[nosep,leftmargin=2em]
\item [A.] Look straight ahead.
\item [B.] Slow after starting into the turn.
\item [C.] Reduce speed coming out of the turn.
\item [D.] Look through the turn to where you want to go.

\end{itemize}
\end{minipage}


\begin{minipage}{\linewidth}
\vspace{1em}\noindent\textbf{36. It is important to search the road for escape routes:}

\begin{itemize}[nosep,leftmargin=2em]
\item [A.] Through alleyways.
\item [B.] At intersections.
\item [C.] When riding a borrowed motorcycle.
\item [D.] On curves.

\end{itemize}
\end{minipage}


\begin{minipage}{\linewidth}
\vspace{1em}\noindent\textbf{37. To make good judgments in traffic, you need to first:}

\begin{itemize}[nosep,leftmargin=2em]
\item [A.] Be able to swerve.
\item [B.] Search ahead.
\item [C.] Watch for pedestrians.
\item [D.] Know how to stop quickly.

\end{itemize}
\end{minipage}


\begin{minipage}{\linewidth}
\vspace{1em}\noindent\textbf{38. The car driver is signaling for a left turn. The rider should:}

\vspace{1ex}
\centering
\includegraphics[width=0.8\linewidth]{ q_38.jpeg }
\captionof*{figure}{ A scene at an intersection where a motorcycle and a car meet at the intersection. The motorcycle is driving straight ahead, while the car is turning in front of motorcycle. The motorcycle has the right-of-way. }

\begin{itemize}[nosep,leftmargin=2em]
\item [A.] Hold speed and center lane position.
\item [B.] Slow and move away from approaching vehicles.
\item [C.] Increase speed and move left.
\item [D.] Increase speed and stay in the center of the lane.

\end{itemize}
\end{minipage}


\begin{minipage}{\linewidth}
\vspace{1em}\noindent\textbf{39. You should make a special point of checking traffic to the rear:}

\begin{itemize}[nosep,leftmargin=2em]
\item [A.] Before starting into a curve.
\item [B.] When you approach an intersection.
\item [C.] When carrying a passenger.
\item [D.] Before slowing.

\end{itemize}
\end{minipage}


\begin{minipage}{\linewidth}
\vspace{1em}\noindent\textbf{40. When preparing to pass, it is usually better to:}

\begin{itemize}[nosep,leftmargin=2em]
\item [A.] Signal, check your left mirror,  look left.
\item [B.] Signal, move into left lane, accelerate.
\item [C.] Signal, check your left mirror.
\item [D.] Signal, Check your left mirror, speed up.

\end{itemize}
\end{minipage}


\begin{minipage}{\linewidth}
\vspace{1em}\noindent\textbf{41. Riding in the center lane position:}

\begin{itemize}[nosep,leftmargin=2em]
\item [A.] Keeps others from seeing you.
\item [B.] Is safest place when approaching intersections.
\item [C.] Keeps others from sharing your lane.
\item [D.] Should be avoided if possible.

\end{itemize}
\end{minipage}


\begin{minipage}{\linewidth}
\vspace{1em}\noindent\textbf{42. To discourage other drivers from sharing your lane, it is usually best to:}

\begin{itemize}[nosep,leftmargin=2em]
\item [A.] Ride in the left portion of your lane.
\item [B.] Ride in the center portion of your lane.
\item [C.] Ride next to another vehicle.
\item [D.] Ride in the right portion of your lane.

\end{itemize}
\end{minipage}


\begin{minipage}{\linewidth}
\vspace{1em}\noindent\textbf{43. Keeping a cushion of space around your motorcycle:}

\begin{itemize}[nosep,leftmargin=2em]
\item [A.] Is only necessary for beginning motorcyclists.
\item [B.] Allows you time and space to react.
\item [C.] Allows you to weave through traffic when you need to.
\item [D.] Makes you more vulnerable.

\end{itemize}
\end{minipage}


\begin{minipage}{\linewidth}
\vspace{1em}\noindent\textbf{44. To spot approaching traffic and maintain a space cushion in the illustration above:}

\vspace{1ex}
\centering
\includegraphics[width=0.8\linewidth]{ q_44.jpeg }
\captionof*{figure}{ Looking ahead at a winding road curve }

\begin{itemize}[nosep,leftmargin=2em]
\item [A.] Ride in the center of the lane.
\item [B.] Move right and lean away from the turn.
\item [C.] Do not lean.
\item [D.] Ride in the left portion of the lane.

\end{itemize}
\end{minipage}


\begin{minipage}{\linewidth}
\vspace{1em}\noindent\textbf{45. You plan to make a right turn after stopping at an intersection. Usually it is best to:}

\begin{itemize}[nosep,leftmargin=2em]
\item [A.] Let the traffic coming from behind pass around you.
\item [B.] Turn from the center portion of the lane.
\item [C.] Turn from the right portion of the lane.
\item [D.] Turn from the left portion of the lane.

\end{itemize}
\end{minipage}


\begin{minipage}{\linewidth}
\vspace{1em}\noindent\textbf{46. Following close to vehicles in front of you:}

\begin{itemize}[nosep,leftmargin=2em]
\item [A.] Helps you see hazards in the road.
\item [B.] Is safer at night.
\item [C.] Protects you in heavy traffic.
\item [D.] Is a factor in crashes caused by motorcyclists.

\end{itemize}
\end{minipage}


\begin{minipage}{\linewidth}
\vspace{1em}\noindent\textbf{47. When following a car, you should try to:}

\begin{itemize}[nosep,leftmargin=2em]
\item [A.] Ride where the driver can see you in the rear-view mirror.
\item [B.] Ride next to another car to be more visible.
\item [C.] Ride in the left side of the lane.
\item [D.] Ride in the right side of the lane.

\end{itemize}
\end{minipage}


\begin{minipage}{\linewidth}
\vspace{1em}\noindent\textbf{48. At intersections, the most common cause of motorcycle/car crashes is:}

\begin{itemize}[nosep,leftmargin=2em]
\item [A.] Drivers entering a rider's right of way.
\item [B.] Riders not seeing drivers enter from side streets.
\item [C.] Drivers tailgating riders.
\item [D.] Riders not yielding to oncoming traffic.

\end{itemize}
\end{minipage}


\begin{minipage}{\linewidth}
\vspace{1em}\noindent\textbf{49. Car/Motorcycle crashes are most common:}

\begin{itemize}[nosep,leftmargin=2em]
\item [A.] At intersections.
\item [B.] At night.
\item [C.] When the rider is speeding.
\item [D.] On freeways.

\end{itemize}
\end{minipage}


\begin{minipage}{\linewidth}
\vspace{1em}\noindent\textbf{50. Making eye contact with a driver at an intersection:}

\begin{itemize}[nosep,leftmargin=2em]
\item [A.] Is a sure sign that the driver will probably yield.
\item [B.] Is important if the driver is pulling out from a side street.
\item [C.] Does not mean that the driver sees you.
\item [D.] Confuses them.

\end{itemize}
\end{minipage}


\begin{minipage}{\linewidth}
\vspace{1em}\noindent\textbf{51. The greatest danger for a rider when passing parked cars is:}

\begin{itemize}[nosep,leftmargin=2em]
\item [A.] A car backing into a parking space.
\item [B.] A car pulling out.
\item [C.] A car door being opened.
\item [D.] Someone stepping from between cars.

\end{itemize}
\end{minipage}


\begin{minipage}{\linewidth}
\vspace{1em}\noindent\textbf{52. In this picture, the light most likely to make the motorcycle noticeable to the car driver is the:}

\vspace{1ex}
\centering
\includegraphics[width=0.8\linewidth]{ q_52.jpeg }
\captionof*{figure}{ Rider merges highway when there is a car in right lane on the highway }

\begin{itemize}[nosep,leftmargin=2em]
\item [A.] Brake light.
\item [B.] Taillight.
\item [C.] Headlight.
\item [D.] Turn signal.

\end{itemize}
\end{minipage}


\begin{minipage}{\linewidth}
\vspace{1em}\noindent\textbf{53. When riding in traffic at night, the best way to locate bumps in the road is to:}

\begin{itemize}[nosep,leftmargin=2em]
\item [A.] Look over the cars ahead.
\item [B.] Put your headlights on high beam.
\item [C.] Read the road signs.
\item [D.] Watch the taillights of the car ahead.

\end{itemize}
\end{minipage}


\begin{minipage}{\linewidth}
\vspace{1em}\noindent\textbf{54. The best way to help others see your motorcycle is to:}

\begin{itemize}[nosep,leftmargin=2em]
\item [A.] Ride in the left portion of the lane.
\item [B.] Wear a white helmet.
\item [C.] Stay in the center portion of the lane.
\item [D.] Ride with your headlight turned on.

\end{itemize}
\end{minipage}


\begin{minipage}{\linewidth}
\vspace{1em}\noindent\textbf{55. Not turning off a turn signal is most likely to:}

\begin{itemize}[nosep,leftmargin=2em]
\item [A.] Encourage other drivers to pull into your path.
\item [B.] Keeps others from noticing your brake light.
\item [C.] Reduces your chances of a crash because other drivers stay clear.
\item [D.] Encourage others to follow you too closely.

\end{itemize}
\end{minipage}


\begin{minipage}{\linewidth}
\vspace{1em}\noindent\textbf{56. Make a special point to use your mirrors:}

\begin{itemize}[nosep,leftmargin=2em]
\item [A.] On long trips.
\item [B.] When riding through an intersection.
\item [C.] When following a vehicle.
\item [D.] Before slowing or stopping.

\end{itemize}
\end{minipage}


\begin{minipage}{\linewidth}
\vspace{1em}\noindent\textbf{57. Before changing lanes to the left:}

\begin{itemize}[nosep,leftmargin=2em]
\item [A.] Check your left mirror; and look to the left.
\item [B.] Check the right mirror, then the left mirror.
\item [C.] Check the left mirror, then the right mirror.
\item [D.] Use your horn and speed up.

\end{itemize}
\end{minipage}


\begin{minipage}{\linewidth}
\vspace{1em}\noindent\textbf{58. The best gear to wear to be seen at night is:}

\begin{itemize}[nosep,leftmargin=2em]
\item [A.] Fluorescent clothing.
\item [B.] A white helmet.
\item [C.] Reflective clothing.
\item [D.] A red jacket.

\end{itemize}
\end{minipage}


\begin{minipage}{\linewidth}
\vspace{1em}\noindent\textbf{59. To see better when following a car at night:}

\begin{itemize}[nosep,leftmargin=2em]
\item [A.] Increase your following distance and use your low beam.
\item [B.] Use the headlights of the vehicle ahead.
\item [C.] Pass when safe and use your high beam.
\item [D.] Drop far back and use your high beam.

\end{itemize}
\end{minipage}


\begin{minipage}{\linewidth}
\vspace{1em}\noindent\textbf{60. When riding at night:}

\begin{itemize}[nosep,leftmargin=2em]
\item [A.] Always ride in the center lane.
\item [B.] Ride with flashers on.
\item [C.] Wear reflective clothing.
\item [D.] Follow closer to the vehicle ahead for safety.

\end{itemize}
\end{minipage}


\begin{minipage}{\linewidth}
\vspace{1em}\noindent\textbf{61. If you lock the front tire when stopping:}

\begin{itemize}[nosep,leftmargin=2em]
\item [A.] Release the rear brake.
\item [B.] Keep the front brake locked until fully stopped.
\item [C.] Keep both brakes locked.
\item [D.] Release and reapply the front brake.

\end{itemize}
\end{minipage}


\begin{minipage}{\linewidth}
\vspace{1em}\noindent\textbf{62. In the picture above, the rear wheel locks while trying to stop quickly. It is usually best to:}

\vspace{1ex}
\centering
\includegraphics[width=0.8\linewidth]{ q_62.jpeg }
\captionof*{figure}{ A street intersection in a suburban neighborhood with a car on the left and a motorcycle on the right aproaching the eachother. }

\begin{itemize}[nosep,leftmargin=2em]
\item [A.] Maintain position and release the rear brake.
\item [B.] Lean the motorcycle and release the rear brake.
\item [C.] Release front brake and keep rear brake locked.
\item [D.] Keep the rear wheel locked until stopped.

\end{itemize}
\end{minipage}


\begin{minipage}{\linewidth}
\vspace{1em}\noindent\textbf{63. When swerving, it is important to:}

\begin{itemize}[nosep,leftmargin=2em]
\item [A.] Always brake before swerving.
\item [B.] Separate braking from swerving.
\item [C.] Always brake after swerving.
\item [D.] Swerve in the direction the hazard is traveling.

\end{itemize}
\end{minipage}


\begin{minipage}{\linewidth}
\vspace{1em}\noindent\textbf{64. To ride over metal bridge gratings:}

\begin{itemize}[nosep,leftmargin=2em]
\item [A.] Ride straight across.
\item [B.] Weave (zig-zag) across.
\item [C.] Cross at a sharp angle without changing lanes.
\item [D.] Speed up.

\end{itemize}
\end{minipage}


\begin{minipage}{\linewidth}
\vspace{1em}\noindent\textbf{65. To ride over an object on the road:}

\begin{itemize}[nosep,leftmargin=2em]
\item [A.] Stay seated until you have crossed it.
\item [B.] Rise slightly off the seat.
\item [C.] Edge across it.
\item [D.] Grip the gas tank with your knees.

\end{itemize}
\end{minipage}


\begin{minipage}{\linewidth}
\vspace{1em}\noindent\textbf{66. When riding over an obstacle, it is usually best to:}

\begin{itemize}[nosep,leftmargin=2em]
\item [A.] Keep a normal seat position.
\item [B.] Rise slightly off the seat.
\item [C.] Grip the gas tank with your knees.
\item [D.] Lean forward as much as possible.

\end{itemize}
\end{minipage}


\begin{minipage}{\linewidth}
\vspace{1em}\noindent\textbf{67. On a paved, two-lane road, sand and gravel are most likely to collect:}

\begin{itemize}[nosep,leftmargin=2em]
\item [A.] In the left wheel track.
\item [B.] Near the center of the road.
\item [C.] Near the sides of the road.
\item [D.] In the center of the lane.

\end{itemize}
\end{minipage}


\begin{minipage}{\linewidth}
\vspace{1em}\noindent\textbf{68. On a wet road, it will usually be most slippery:}

\begin{itemize}[nosep,leftmargin=2em]
\item [A.] On loose gravel.
\item [B.] In the left wheel track.
\item [C.] In the right wheel track.
\item [D.] In the center portion of the lane.

\end{itemize}
\end{minipage}


\begin{minipage}{\linewidth}
\vspace{1em}\noindent\textbf{69. To handle a slippery surface, it is best to:}

\begin{itemize}[nosep,leftmargin=2em]
\item [A.] Use the rear brake lightly while crossing.
\item [B.] Slow down before reaching it.
\item [C.] Gradually accelerate.
\item [D.] Downshift when you reach it.

\end{itemize}
\end{minipage}


\begin{minipage}{\linewidth}
\vspace{1em}\noindent\textbf{70. To stop on wet pavement:}

\begin{itemize}[nosep,leftmargin=2em]
\item [A.] Apply the front brake only if the wheel begins to slide.
\item [B.] Coast to a stop.
\item [C.] Apply the rear brake only.
\item [D.] Apply the front and rear brakes together.

\end{itemize}
\end{minipage}


\begin{minipage}{\linewidth}
\vspace{1em}\noindent\textbf{71. When riding where sand and gravel have collected on paved roads, you should:}

\begin{itemize}[nosep,leftmargin=2em]
\item [A.] Downshift and use more throttle.
\item [B.] Pull-in the clutch.
\item [C.] Only use the rear brake to slow down.
\item [D.] Avoid sudden changes in speed or direction.

\end{itemize}
\end{minipage}


\begin{minipage}{\linewidth}
\vspace{1em}\noindent\textbf{72. If the rear tire goes flat:}

\begin{itemize}[nosep,leftmargin=2em]
\item [A.] You will not be able to maintain control.
\item [B.] You will lose power to the rear wheel.
\item [C.] The back of the motorcycle will jerk from side to side.
\item [D.] The steering will feel heavy.

\end{itemize}
\end{minipage}


\begin{minipage}{\linewidth}
\vspace{1em}\noindent\textbf{73. If you have a flat tire while riding, hold the handle grips firmly and:}

\begin{itemize}[nosep,leftmargin=2em]
\item [A.] Gently apply both brakes.
\item [B.] Use the brake of the good tire or avoid braking.
\item [C.] Continue to the nearest repair facility.
\item [D.] Use only the brake that controls the flat tire.

\end{itemize}
\end{minipage}


\begin{minipage}{\linewidth}
\vspace{1em}\noindent\textbf{74. When the front tire goes flat while riding:}

\begin{itemize}[nosep,leftmargin=2em]
\item [A.] Apply the front brake.
\item [B.] You will not be able to control the motorcycle.
\item [C.] Apply rear brake and steer to the edge of the road.
\item [D.] Steer to the right and apply both brakes.

\end{itemize}
\end{minipage}


\begin{minipage}{\linewidth}
\vspace{1em}\noindent\textbf{75. When the front tire goes flat:}

\begin{itemize}[nosep,leftmargin=2em]
\item [A.] The steering will feel heavy.
\item [B.] It has no effect on handling.
\item [C.] The back of the motorcycle will jerk from side to side.
\item [D.] You will lose power to the brakes.

\end{itemize}
\end{minipage}


\begin{minipage}{\linewidth}
\vspace{1em}\noindent\textbf{76. If your throttle sticks:}

\begin{itemize}[nosep,leftmargin=2em]
\item [A.] Squeeze the clutch and use engine cut-off switch.
\item [B.] Shift to neutral and apply both brakes.
\item [C.] Downshift and apply both brakes.
\item [D.] Apply the front brake.

\end{itemize}
\end{minipage}


\begin{minipage}{\linewidth}
\vspace{1em}\noindent\textbf{77. In the picture above, it is usually best to:}

\vspace{1ex}
\centering
\includegraphics[width=0.8\linewidth]{ q_77.jpeg }
\captionof*{figure}{ Illustration of the lane with three numbered portions: 1 - left, 2 - center, 3 - right. }

\begin{itemize}[nosep,leftmargin=2em]
\item [A.] Speed up to get by quickly.
\item [B.] Ride in position 3.
\item [C.] Ride in position 1.
\item [D.] Ride in position 2 or 3.

\end{itemize}
\end{minipage}


\begin{minipage}{\linewidth}
\vspace{1em}\noindent\textbf{78. In this picture, A, B, and C are riding as a group. To pass, rider B should:}

\vspace{1ex}
\centering
\includegraphics[width=0.8\linewidth]{ q_78.jpeg }
\captionof*{figure}{ Illustration of bike formation passing a car. A passed, B is passing, C is waiting. }

\begin{itemize}[nosep,leftmargin=2em]
\item [A.] Complete the pass in front of A.
\item [B.] Pass with C.
\item [C.] Return to the lane and ride next to A.
\item [D.] Return to the right portion of the lane behind A.

\end{itemize}
\end{minipage}


\begin{minipage}{\linewidth}
\vspace{1em}\noindent\textbf{79. To create more space in the situation pictured above, you should:}

\vspace{1ex}
\centering
\includegraphics[width=0.8\linewidth]{ q_79.jpeg }
\captionof*{figure}{ Motorcycle driving in right lane of a highway, a car mergin to the right lane infront of  motorcycle. }

\begin{itemize}[nosep,leftmargin=2em]
\item [A.] Ride in the right portion of the lane.
\item [B.] Adjust speed and stay in the center of the lane.
\item [C.] Move to the left lane.
\item [D.] Pass the car on your left.

\end{itemize}
\end{minipage}


\begin{minipage}{\linewidth}
\vspace{1em}\noindent\textbf{80. Most motorcycle/car crashes:}

\begin{itemize}[nosep,leftmargin=2em]
\item [A.] Occur because a car driver pulls into the path of the motorcycle.
\item [B.] Result in minor injuries.
\item [C.] Occur because a motorcyclist turns left in front of a vehicle.
\item [D.] Occur because of speeding.

\end{itemize}
\end{minipage}


\begin{minipage}{\linewidth}
\vspace{1em}\noindent\textbf{81. You are preparing to move to the left lane, as pictured above. To spot cars passing beside you, it is best to:}

\vspace{1ex}
\centering
\includegraphics[width=0.8\linewidth]{ q_81.jpeg }
\captionof*{figure}{ Motorcycle in center of a middle lane. Prepearing to move to left lane. }

\begin{itemize}[nosep,leftmargin=2em]
\item [A.] Slow and look in the left mirror.
\item [B.] Look in the left mirror.
\item [C.] Speed up and look in the left mirror.
\item [D.] Turn your head to the left.

\end{itemize}
\end{minipage}


\begin{minipage}{\linewidth}
\vspace{1em}\noindent\textbf{82. To adjust for traffic and avoid debris in the curve pictured, ride:}

\vspace{1ex}
\centering
\includegraphics[width=0.8\linewidth]{ q_82.jpeg }
\captionof*{figure}{ Undivided highway curve where right side has numbered portions: 1 - left portion, 2 - center portion, 3 - right portion }

\begin{itemize}[nosep,leftmargin=2em]
\item [A.] In position 3.
\item [B.] In position 2.
\item [C.] Without leaning.
\item [D.] In position 1.

\end{itemize}
\end{minipage}


\begin{minipage}{\linewidth}
\vspace{1em}\noindent\textbf{83. To cross tracks that run parallel to your course:}

\begin{itemize}[nosep,leftmargin=2em]
\item [A.] Edge gradually across.
\item [B.] Move into the opposing lane to cross at a 90-degree angle.
\item [C.] Wait for the tracks to end.
\item [D.] Move away, then turn and cross at an angle of at least 45 degrees.

\end{itemize}
\end{minipage}


\begin{minipage}{\linewidth}
\vspace{1em}\noindent\textbf{84. In this picture, the rider in the most dangerous position is:}

\vspace{1ex}
\centering
\includegraphics[width=0.8\linewidth]{ q_84.jpeg }
\captionof*{figure}{ Illustration of bike formation passing a car. A passed, B is passing, C is waiting. }

\begin{itemize}[nosep,leftmargin=2em]
\item [A.] Rider B.
\item [B.] Rider C.
\item [C.] Rider A.
\item [D.] None.

\end{itemize}
\end{minipage}


